\documentclass[a4paper]{article}

%% Language and font encodings
\usepackage[portuges]{babel}
\usepackage{fontspec}

%% Sets page size and margins
\usepackage[a4paper,top=3cm,bottom=2cm,left=3cm,right=3cm,marginparwidth=1.75cm]{geometry}

%% Useful packages
\usepackage{amsmath,amsthm,amssymb,amsfonts}
\usepackage{graphicx}
\usepackage[colorinlistoftodos]{todonotes}
\usepackage[colorlinks=true, allcolors=blue]{hyperref}

\newcommand{\R}{\mathbb{R}}
\newcommand{\N}{\mathbb{N}}
\newcommand{\Z}{\mathbb{Z}}
\providecommand{\C}{\mathbb{C}}

\theoremstyle{definition}
\newtheorem{defin}{Definição}

\theoremstyle{plain}
\newtheorem{theorem}[defin]{Teorema}
\newtheorem{corollary}[defin]{Corolário}



\title{Seminário II 2019/20 - Ficha de trabalho N }

\author{Nome }

\date{dd/mm/yyyy}

\begin{document}
\maketitle

%\begin{abstract}
%Your abstract.
%\end{abstract}

\section{Introdução}

%\emph{\small Incluir aqui uma breve descrição do tema}


\begin{defin}
Uma aplicação $$f: A \rightarrow B$$ em que $A$ é o domínio de $f$ e $B$ é o conjunto de chegada de $f$, diz-se função real de variável real se $A,B \subseteq \R$. 
\end{defin}

\begin{theorem}[Álgebra]
Uma função real de variável real é invertível se e só se for bijectiva.
\end{theorem}

\begin{corollary}\label{estritcresc}
Se $f$ é sobrejectiva e estritamente crescente no seu dom\'inio então $f$ é invertível. \end{corollary}
\section{Trabalho realizado}


%Podem inserir gráficos, usando \includegraphics: por exemplo,
%\includegraphics[width=0.5\linewidth]{cloud}
%
% Veremos na aula como fazer upload de imagens para usar aqui
%Etapas de realização do trabalho.
\begin{enumerate}
\item Determinámos analiticamente o dominio e o contradomínio da fun\c c\~ao $\displaystyle f(x)=\frac{1}{x+1}$.
\item Usando o Geogebra, realizámos os seguintes passos\begin{enumerate}
\item usando a linha de \emph{input}, construímos o gráfico da função $f$.
\item construímos um ponto $A$, de coordenadas $(-1,a)$
\item observámos que o gráfico da função $f$ não contém o ponto $A$.
\item construímos o gráfico da função $$h_a(x)=\left\{\begin{array}{lr}f(x)& se\ x\neq -1\\a& se\  x=-1\end{array}\right.$$
\item verificámos que a função $h_a(x)$ é bijectiva se e só $a=0$.
\item tracámos o gráfico da recta y=x.
\item verificámos que o gráfico de $h_0^{-1}(x)$ se pode obter por reflexão em relação à recta $y=x$.
\centerline{\includegraphics[width=0.5\linewidth]{seminario_funcoes}}

\end{enumerate}
\end{enumerate}



\section{Conclusões}
Concluímos que o recíproco do corolário \ref{estritcresc} não é verdadeiro visto que a função $h_0(x)$ é invertível e não é estritamente crescente.

%expor aqui as conclusões que retiraram do trabalho

\end{document}